% ---
% Comandos auxiliares
% ---

% Lista de referências: indica as páginas em que cada obra foi citada
\renewcommand{\backrefpagesname}{Citado na(s) página(s):~}
\renewcommand{\backref}{}
\renewcommand*{\backrefalt}[4]{
	\ifcase #1 %
		Nenhuma citação no texto.%
	\or
		Citado na página #2.%
	\else
		Citado #1 vezes nas páginas #2.%
	\fi}%

% Endereços do projeto (premissa 12 do roteiro: link direto ao repositório)
\newcommand{\urlrepositorio}{\url{https://github.com/FCTE-Qualidade-de-Software-1/2026-2_T02_ADELE_GOLDSTINE}}
\newcommand{\urlgitpage}{\url{https://fcte-qualidade-de-software-1.github.io/2026-2_T02_ADELE_GOLDSTINE/}}

% Colunas de tabela com quebra de linha e alinhamento à esquerda
\newcolumntype{L}[1]{>{\raggedright\arraybackslash}p{#1}}
\newcolumntype{C}[1]{>{\centering\arraybackslash}p{#1}}
\newcolumntype{Y}{>{\raggedright\arraybackslash}X}

% Marcação de pendências que a equipe precisa preencher antes da entrega.
% Procure por \pendente no código-fonte; nenhuma deve restar na versão final.
\newcommand{\pendente}[1]{\textcolor{red}{#1}}
